\documentclass[11pt,a4paper]{article}
\usepackage[utf8]{inputenc}
\usepackage[T1]{fontenc}
\usepackage[catalan]{babel}
\usepackage{lmodern}
\usepackage[sfdefault,lining]{FiraSans}
\usepackage[a4paper,top=2.3cm,bottom=2.2cm,left=1.8cm,right=1.8cm,headheight=1.3cm,headsep=0.4cm,footskip=1cm]{geometry}
\usepackage{amsmath,amssymb}
\usepackage{icomma}
\usepackage{graphicx}
\usepackage[table]{xcolor}
\usepackage{tikz}
\usepackage{fontawesome5}
\usepackage{ccicons}
\usepackage{fancyhdr}
\usepackage{lastpage}
\usepackage{tasks}
\usepackage{needspace}
\usepackage[most]{tcolorbox}
\usepackage{array}
\usepackage[hidelinks,pdftitle={Fitxa d'exercicis de la unitat 4: Fraccions. 1r d'ESO},pdfsubject={Matemàtiques 1r ESO},pdfkeywords={fraccions, fracció d'un nombre, percentatges, fraccions equivalents, comú denominador, comparar fraccions, operacions amb fraccions, problemes}]{hyperref}

\definecolor{navy}{HTML}{1B2748}
\definecolor{accent}{HTML}{E8603A}
\definecolor{teal}{HTML}{1F9E90}
\definecolor{soft}{HTML}{EEF1F8}
\definecolor{muted}{HTML}{5D6685}

\setlength{\parindent}{0pt}
\setlength{\parskip}{3pt}
\settasks{label=\alph*), label-format=\bfseries\color{navy}, label-width=1.3em, item-indent=1.8em, column-sep=0.6em, after-item-skip=5pt, before-skip=3pt, after-skip=4pt}

% ---------- capçalera i peu ----------
\newcommand{\llicencia}{\href{https://creativecommons.org/licenses/by-nc-sa/4.0/deed.ca}{\ccbyncsa}\ \ Llicència Creative Commons Reconeixement-NoComercial-CompartirIgual 4.0 Internacional (CC BY-NC-SA 4.0)}
\pagestyle{fancy}
\fancyhf{}
\fancyhead[L]{\raisebox{-0.35\height}{\includegraphics[height=0.95cm]{logo-gva.png}}}
\fancyhead[C]{\small\color{muted}Fitxa d'exercicis · Unitat 4: Fraccions · 1r d'ESO}
\fancyhead[R]{\raisebox{-0.35\height}{\includegraphics[height=0.95cm]{logo-vermellar.png}}}
\fancyfoot[L]{\scriptsize\color{muted}\llicencia}
\fancyfoot[R]{\small\color{muted}\thepage}
\renewcommand{\headrulewidth}{0.4pt}
\renewcommand{\headrule}{\color{muted!40}\hrule height \headrulewidth}
\fancypagestyle{primera}{\fancyhf{}\renewcommand{\headrulewidth}{0pt}
  \fancyfoot[L]{\scriptsize\color{muted}\llicencia}
  \fancyfoot[R]{\small\color{muted}\thepage}}

% ---------- nivells i exercicis ----------
\newcommand{\estrelles}[1]{{\color{accent}\ifcase#1\or\faStar\faStar[regular]\faStar[regular]\or\faStar\faStar\faStar[regular]\or\faStar\faStar\faStar\fi}}
\newcounter{exer}
\newcommand{\ex}[1]{\par\needspace{4\baselineskip}\medskip\stepcounter{exer}\noindent{\color{navy}\bfseries\theexer.}\ \ {\small\estrelles{#1}}\ \ }

\newcommand{\seccio}[2]{\par\needspace{10\baselineskip}\bigskip
  \begin{tcolorbox}[colback=navy,colframe=navy,coltext=white,arc=2mm,boxrule=0pt,left=3mm,right=3mm,top=1.2mm,bottom=1.2mm]
  {\large\bfseries #1\quad #2}
  \end{tcolorbox}\vspace{-2pt}}

\newcolumntype{C}[1]{>{\centering\arraybackslash}p{#1}}

% quadre de recordatori (només a la fitxa d'exercicis)
\newcommand{\recorda}[1]{\par\smallskip\begin{tcolorbox}[colback=accent!7,colframe=accent!55,arc=2mm,boxrule=0.5pt,left=2.5mm,right=2.5mm,top=1mm,bottom=1mm]\small{\color{accent}\faLightbulb[regular]}\ \textbf{Recorda:}\ #1\end{tcolorbox}}
\newcommand{\res}[1]{#1}

\begin{document}
\thispagestyle{primera}

\begin{minipage}[c]{0.25\textwidth}\includegraphics[width=\linewidth]{logo-gva.png}\end{minipage}%
\begin{minipage}[c]{0.5\textwidth}\centering
  {\color{navy}\Large\bfseries Fitxa d'exercicis de la unitat 4:\\[2pt] Fraccions}\\[5pt]
  {\color{accent}\large\bfseries 1r d'ESO}
\end{minipage}%
\begin{minipage}[c]{0.25\textwidth}\hfill\includegraphics[width=0.85\linewidth]{logo-vermellar.png}\end{minipage}

\vspace{6pt}
{\color{navy}\hrule height 1.2pt}
\vspace{10pt}
Nom i cognoms: \rule{8.4cm}{0.4pt}\hfill Grup: \rule{1.6cm}{0.4pt}\hfill Data: \rule{2.2cm}{0.4pt}
\vspace{8pt}

\begin{tcolorbox}[colback=soft,colframe=soft,arc=2mm,left=3mm,right=3mm,top=1.5mm,bottom=1.5mm]
\small Exercicis graduats per nivells de dificultat:\quad \estrelles{1}\ bàsic \qquad \estrelles{2}\ mitjà \qquad \estrelles{3}\ avançat.\\[2pt] Resol-los al quadern mostrant tots els passos i dona els resultats en forma de fracció irreductible. Les figures i les rectes es poden completar directament en la fitxa. Fes els càlculs sense calculadora.
\end{tcolorbox}

% =====================================================
\seccio{1}{El concepte de fracció}

\ex{1} Quina fracció representa la part pintada de cada figura?\begin{tasks}(3)
  \task \begin{tikzpicture}[scale=0.42,baseline=(current bounding box.center)]\fill[accent!75] (0,1) rectangle ++(1,1);\fill[accent!75] (1,1) rectangle ++(1,1);\fill[accent!75] (2,1) rectangle ++(1,1);\fill[accent!75] (3,1) rectangle ++(1,1);\draw[thick] (0,0) grid (5,2);\end{tikzpicture}
  \task \begin{tikzpicture}[scale=0.42,baseline=(current bounding box.center)]\fill[accent!75] (0,1) rectangle ++(1,1);\fill[accent!75] (1,1) rectangle ++(1,1);\fill[accent!75] (0,0) rectangle ++(1,1);\draw[thick] (0,0) grid (2,2);\end{tikzpicture}
  \task \begin{tikzpicture}[baseline=(current bounding box.center)]\fill[accent!75] (0,0) -- (90.00:0.75) arc[start angle=90.00,end angle=38.57,radius=0.75] -- cycle;\draw[thick] (0,0) circle (0.75);\draw[thick] (0,0) -- (90.00:0.75);\draw[thick] (0,0) -- (38.57:0.75);\draw[thick] (0,0) -- (-12.86:0.75);\draw[thick] (0,0) -- (-64.29:0.75);\draw[thick] (0,0) -- (-115.71:0.75);\draw[thick] (0,0) -- (-167.14:0.75);\draw[thick] (0,0) -- (-218.57:0.75);\end{tikzpicture}
  \task \begin{tikzpicture}[scale=0.42,baseline=(current bounding box.center)]\fill[accent!75] (0,1) rectangle ++(1,1);\fill[accent!75] (1,1) rectangle ++(1,1);\draw[thick] (0,0) grid (2,2);\end{tikzpicture}
  \task \begin{tikzpicture}[baseline=(current bounding box.center)]\fill[accent!75] (0,0) -- (90.00:0.75) arc[start angle=90.00,end angle=38.57,radius=0.75] -- cycle;\fill[accent!75] (0,0) -- (38.57:0.75) arc[start angle=38.57,end angle=-12.86,radius=0.75] -- cycle;\fill[accent!75] (0,0) -- (-12.86:0.75) arc[start angle=-12.86,end angle=-64.29,radius=0.75] -- cycle;\fill[accent!75] (0,0) -- (-64.29:0.75) arc[start angle=-64.29,end angle=-115.71,radius=0.75] -- cycle;\fill[accent!75] (0,0) -- (-115.71:0.75) arc[start angle=-115.71,end angle=-167.14,radius=0.75] -- cycle;\draw[thick] (0,0) circle (0.75);\draw[thick] (0,0) -- (90.00:0.75);\draw[thick] (0,0) -- (38.57:0.75);\draw[thick] (0,0) -- (-12.86:0.75);\draw[thick] (0,0) -- (-64.29:0.75);\draw[thick] (0,0) -- (-115.71:0.75);\draw[thick] (0,0) -- (-167.14:0.75);\draw[thick] (0,0) -- (-218.57:0.75);\end{tikzpicture}
  \task \begin{tikzpicture}[scale=0.42,baseline=(current bounding box.center)]\fill[accent!75] (0,1) rectangle ++(1,1);\fill[accent!75] (1,1) rectangle ++(1,1);\fill[accent!75] (2,1) rectangle ++(1,1);\fill[accent!75] (0,0) rectangle ++(1,1);\draw[thick] (0,0) grid (3,2);\end{tikzpicture}
\end{tasks}\recorda{El \textbf{denominador} indica en quantes parts \textbf{iguals} dividim la unitat; el \textbf{numerador}, quantes n'agafem.}

\ex{1} Indica el numerador i el denominador de cada fracció i escriu com es llegeix:\begin{tasks}(4)
  \task $\dfrac{3}{5}$
  \task $\dfrac{4}{6}$
  \task $\dfrac{12}{15}$
  \task $\dfrac{4}{3}$
\end{tasks}

\ex{1} Pinta en cada figura la fracció indicada:\begin{tasks}(3)
  \task $\dfrac{3}{5}$\ \ \begin{tikzpicture}[baseline=(current bounding box.center)]\draw (0.700,0)--(0.700,0.55);\draw (1.400,0)--(1.400,0.55);\draw (2.100,0)--(2.100,0.55);\draw (2.800,0)--(2.800,0.55);\draw[thick] (0,0) rectangle (3.5,0.55);\end{tikzpicture}
  \task $\dfrac{5}{8}$\ \ \begin{tikzpicture}[scale=0.42,baseline=(current bounding box.center)]\draw[thick] (0,0) grid (4,2);\end{tikzpicture}
  \task $\dfrac{2}{3}$\ \ \begin{tikzpicture}[baseline=(current bounding box.center)]\draw[thick] (0,0) circle (0.65);\draw[thick] (0,0) -- (90.00:0.65);\draw[thick] (0,0) -- (-30.00:0.65);\draw[thick] (0,0) -- (-150.00:0.65);\end{tikzpicture}
\end{tasks}

\ex{1} Representa en la recta numèrica les fraccions $\dfrac{6}{10}$ i $\dfrac{7}{4}$.\par\smallskip{\centering \begin{tikzpicture}[baseline=0]\draw[-latex,thick] (-0.3,0)--(9.50,0);\draw[thick] (0.00,0.18)--(0.00,-0.18) node[below=1pt,font=\small]{$0$};\draw[thick] (4.50,0.18)--(4.50,-0.18) node[below=1pt,font=\small]{$1$};\draw[thick] (9.00,0.18)--(9.00,-0.18) node[below=1pt,font=\small]{$2$};\end{tikzpicture}\par}

\ex{2} Classifica les fraccions en pròpies, impròpies o iguals a la unitat. Escriu les impròpies com a unitats senceres més una fracció pròpia.\begin{tasks}(6)
  \task $\dfrac{7}{3}$
  \task $\dfrac{5}{8}$
  \task $\dfrac{9}{9}$
  \task $\dfrac{13}{4}$
  \task $\dfrac{2}{11}$
  \task $\dfrac{20}{6}$
\end{tasks}

\ex{2} Quina fracció representa?\begin{tasks}(2)
  \task 15 minuts d'una hora
  \task 4 mesos d'un any
  \task 5 dies d'una setmana
  \task 300 grams d'1 kg
\end{tasks}

\ex{3} Escriu totes les fraccions de denominador $9$ que són impròpies però menors que $2$. Quines fraccions de denominador $6$ hi ha entre $2$ i $3$ (sense comptar-los)?


% =====================================================
\seccio{2}{Fracció d'un nombre}

\ex{1} Calcula:\begin{tasks}(3)
  \task $\dfrac{5}{6}$ de $48$
  \task $\dfrac{4}{7}$ de $490$
  \task $\dfrac{3}{5}$ de $200$
  \task $\dfrac{7}{8}$ de $32$
  \task $\dfrac{2}{5}$ de $60$
  \task $\dfrac{3}{6}$ de $180$
\end{tasks}\recorda{$\dfrac{3}{5}$ de $40$: dividim entre el denominador i multipliquem pel numerador: \ $(40:5)\cdot3=24$.}

\ex{1} En una classe de 30 alumnes, $\dfrac{2}{5}$ són xics. Quants xics i quantes xiques hi ha?

\ex{2} Un dipòsit de $2400$ litres està ple fins als $\dfrac{3}{8}$. Quants litres té? Quants en falten per a omplir-lo?

\ex{2} Calcula el nombre:\begin{tasks}(2)
  \task Els $\dfrac{3}{4}$ d'un nombre són $36$.
  \task Els $\dfrac{2}{7}$ d'un nombre són $18$.
\end{tasks}

\ex{2} En Pau s'ha gastat $\dfrac{3}{5}$ de la paga i encara li queden $14$\,€. Quants diners tenia?

\ex{3} Calcula $\dfrac{2}{3}$ de $\dfrac{3}{4}$ de $360$.

\ex{3} Van d'excursió 84 alumnes: $\dfrac{3}{7}$ són de 1r, $\dfrac{1}{4}$ de 2n i la resta de 3r. Quants n'hi ha de cada curs?


% =====================================================
\seccio{3}{Percentatges i fraccions}

\ex{1} Calcula:\begin{tasks}(3)
  \task $20\,\%$ de $180$
  \task $45\,\%$ de $80$
  \task $12\,\%$ de $350$
  \task $85\,\%$ de $1300$
  \task $50\,\%$ de $64$
  \task $25\,\%$ de $120$
\end{tasks}\recorda{$p\,\%=\dfrac{p}{100}$. \ El $p\,\%$ de $N$ és $\dfrac{p\cdot N}{100}$.}

\ex{1} Escriu cada percentatge com a fracció irreductible:\begin{tasks}(4)
  \task $50\,\%$
  \task $25\,\%$
  \task $75\,\%$
  \task $20\,\%$
  \task $10\,\%$
  \task $40\,\%$
  \task $5\,\%$
  \task $35\,\%$
\end{tasks}

\ex{1} Expressa en percentatge:\begin{tasks}(6)
  \task $\dfrac{1}{2}$
  \task $\dfrac{3}{4}$
  \task $\dfrac{2}{5}$
  \task $\dfrac{7}{10}$
  \task $\dfrac{9}{20}$
  \task $\dfrac{3}{25}$
\end{tasks}

\ex{2} Una jaqueta costa $80$\,€ i la rebaixen un $25\,\%$. Quants euros ens descompten? Quant paguem?

\ex{2} En un institut de 650 alumnes, el $52\,\%$ són xiques. Quantes xiques i quants xics hi ha?

\ex{3} En una província de $378\,000$ persones, el $7\,\%$ són del Marroc, el $25\,\%$ de l'est d'Europa, el $12\,\%$ d'Alemanya i la resta són d'ací. Quin percentatge són d'ací? Quantes persones hi ha de cada grup?

\ex{3} En una classe de 25 alumnes, 8 han suspés l'examen. Quin percentatge ha aprovat?


% =====================================================
\seccio{4}{Fraccions equivalents}

\ex{1} Escriu, per amplificació, quatre fraccions equivalents a cadascuna:\begin{tasks}(6)
  \task $\dfrac{1}{4}$
  \task $\dfrac{2}{5}$
  \task $\dfrac{5}{15}$
  \task $\dfrac{7}{13}$
  \task $\dfrac{11}{17}$
  \task $\dfrac{1}{10}$
\end{tasks}\recorda{\textbf{Amplificar}: multiplicar numerador i denominador pel mateix nombre. \textbf{Simplificar}: dividir-los entre un divisor comú. La fracció \textbf{irreductible} s'obté dividint entre el MCD.}

\ex{1} Troba dues fraccions simplificades de cadascuna:\begin{tasks}(5)
  \task $\dfrac{4}{12}$
  \task $\dfrac{12}{20}$
  \task $\dfrac{66}{84}$
  \task $\dfrac{15}{45}$
  \task $\dfrac{50}{100}$
\end{tasks}

\ex{1} Són irreductibles?\begin{tasks}(3)
  \task $\dfrac{25}{35}$
  \task $\dfrac{5}{7}$
  \task $\dfrac{98}{49}$
\end{tasks}

\ex{2} Simplifica fins a arribar a la fracció irreductible:\begin{tasks}(5)
  \task $\dfrac{4}{6}$
  \task $\dfrac{4}{20}$
  \task $\dfrac{3}{12}$
  \task $\dfrac{2}{6}$
  \task $\dfrac{35}{40}$
  \task $\dfrac{75}{100}$
  \task $\dfrac{4}{16}$
  \task $\dfrac{14}{21}$
  \task $\dfrac{10}{15}$
  \task $\dfrac{25}{45}$
  \task $\dfrac{13}{52}$
  \task $\dfrac{33}{34}$
  \task $\dfrac{705}{120}$
  \task $\dfrac{170}{105}$
  \task $\dfrac{256}{512}$
\end{tasks}

\ex{2} Calcula la fracció irreductible de cadascuna i indica quines són equivalents entre si:\begin{tasks}(3)
  \task $\dfrac{54}{24}$
  \task $\dfrac{40}{88}$
  \task $\dfrac{81}{36}$
  \task $\dfrac{72}{32}$
  \task $\dfrac{75}{185}$
  \task $\dfrac{50}{110}$
\end{tasks}

\ex{3} Troba una fracció equivalent a $\dfrac{15}{24}$ que tinga denominador $16$, i una d'equivalent a $\dfrac{5}{8}$ que tinga numerador $80$.

\ex{3} Escriu totes les fraccions equivalents a $\dfrac{2}{3}$ que tenen el denominador entre $20$ i $30$.


% =====================================================
\seccio{5}{Comprovar si dues fraccions són equivalents}

\ex{1} Comprova si són equivalents:\begin{tasks}(3)
  \task $\dfrac{8}{32}$ i $\dfrac{3}{12}$
  \task $\dfrac{6}{15}$ i $\dfrac{26}{34}$
  \task $\dfrac{7}{23}$ i $\dfrac{14}{46}$
  \task $\dfrac{1}{2}$ i $\dfrac{15}{30}$
  \task $\dfrac{68}{52}$ i $\dfrac{17}{13}$
  \task $\dfrac{3}{29}$ i $\dfrac{33}{67}$
\end{tasks}\recorda{$\dfrac{a}{b}=\dfrac{c}{d}$ \ si i només si \ $a\cdot d=b\cdot c$ \ (productes creuats).}

\ex{1} Calcula el valor de $x$ perquè les fraccions siguen equivalents:\begin{tasks}(5)
  \task $\dfrac{3}{14}=\dfrac{x}{42}$
  \task $\dfrac{6}{8}=\dfrac{24}{x}$
  \task $\dfrac{x}{5}=\dfrac{12}{20}$
  \task $\dfrac{14}{3}=\dfrac{x}{21}$
  \task $\dfrac{5}{11}=\dfrac{15}{x}$
\end{tasks}

\ex{2} Completa perquè siguen equivalents: \quad $\dfrac{2}{5}=\dfrac{\square}{20}=\dfrac{12}{\square}=\dfrac{\square}{45}$

\ex{2} Quina d'aquestes fraccions no és equivalent a les altres? \quad $\dfrac{6}{9},\ \dfrac{10}{15},\ \dfrac{14}{20},\ \dfrac{16}{24}$

\ex{3} Troba $x$ i $y$ perquè \ $\dfrac{x}{12}=\dfrac{15}{20}=\dfrac{9}{y}$.

\ex{3} Troba una fracció equivalent a $\dfrac{3}{5}$ que tinga el denominador $24$ unitats més gran que el numerador.


% =====================================================
\seccio{6}{Reducció a comú denominador}

\ex{1} Redueix a comú denominador:\begin{tasks}(4)
  \task $\dfrac{1}{5},\ \dfrac{3}{7}$
  \task $\dfrac{2}{9},\ \dfrac{1}{6}$
  \task $\dfrac{3}{4},\ \dfrac{3}{8}$
  \task $\dfrac{2}{7},\ \dfrac{1}{4}$
  \task $\dfrac{5}{21},\ \dfrac{3}{2}$
  \task $\dfrac{3}{5},\ \dfrac{1}{10}$
  \task $\dfrac{1}{4},\ \dfrac{1}{6}$
  \task $\dfrac{3}{7},\ 5$
\end{tasks}\recorda{Nou denominador: el \textbf{mcm} dels denominadors. Per a cada fracció, dividim el mcm entre el seu denominador i multipliquem el numerador pel resultat.}

\ex{2} Redueix a comú denominador:\begin{tasks}(4)
  \task $\dfrac{7}{24},\ \dfrac{3}{16}$
  \task $\dfrac{1}{2},\ \dfrac{2}{3},\ \dfrac{5}{8}$
  \task $\dfrac{7}{2},\ \dfrac{6}{15},\ \dfrac{3}{10}$
  \task $\dfrac{5}{8},\ \dfrac{1}{2},\ \dfrac{7}{12}$
  \task $\dfrac{4}{7},\ \dfrac{1}{3},\ \dfrac{5}{14}$
  \task $\dfrac{3}{18},\ \dfrac{4}{30},\ \dfrac{7}{36}$
  \task $\dfrac{6}{7},\ \dfrac{2}{14},\ \dfrac{12}{28}$
  \task $\dfrac{5}{3},\ \dfrac{5}{12},\ \dfrac{3}{4},\ \dfrac{9}{20}$
\end{tasks}

\ex{2} Escriu les fraccions $\dfrac{5}{6},\ \dfrac{7}{12},\ \dfrac{4}{9}$ amb denominador $36$.

\ex{3} Simplifica primer i després redueix a comú denominador amb el denominador més xicotet possible: \ $\dfrac{6}{9},\ \dfrac{10}{15},\ \dfrac{9}{12}$.


% =====================================================
\seccio{7}{Comparar i ordenar fraccions}

\ex{1} Escriu una fracció més gran i una de més xicoteta que la donada:\begin{tasks}(3)
  \task $\dfrac{5}{7}$
  \task $\dfrac{4}{9}$
  \task $\dfrac{3}{13}$
\end{tasks}

\ex{1} Compara (escriu $<$, $>$ o $=$):\begin{tasks}(4)
  \task $\dfrac{9}{5}\ \square\ \dfrac{9}{8}$
  \task $\dfrac{6}{11}\ \square\ \dfrac{10}{11}$
  \task $\dfrac{3}{7}\ \square\ \dfrac{3}{4}$
  \task $\dfrac{4}{8}\ \square\ \dfrac{2}{16}$
  \task $\dfrac{2}{5}\ \square\ \dfrac{5}{6}$
  \task $\dfrac{1}{3}\ \square\ \dfrac{3}{7}$
  \task $\dfrac{8}{9}\ \square\ \dfrac{11}{12}$
  \task $\dfrac{3}{4}\ \square\ \dfrac{24}{32}$
\end{tasks}\recorda{Mateix denominador: més gran el de numerador més gran. Mateix numerador: més gran el de denominador més xicotet. Si no, redueix a comú denominador.}

\ex{2} Ordena de menor a major:\begin{tasks}(3)
  \task $\dfrac{3}{5},\ \dfrac{2}{7},\ \dfrac{11}{2}$
  \task $\dfrac{4}{9},\ \dfrac{5}{6},\ \dfrac{2}{3}$
  \task $\dfrac{5}{4},\ \dfrac{2}{6},\ \dfrac{1}{10}$
  \task $\dfrac{18}{3},\ \dfrac{14}{7},\ \dfrac{23}{21}$
  \task $\dfrac{2}{12},\ \dfrac{4}{15},\ \dfrac{3}{24}$
  \task $\dfrac{12}{14},\ \dfrac{6}{7},\ \dfrac{3}{28}$
\end{tasks}

\ex{2} Ordena de menor a major:\begin{tasks}(3)
  \task $\dfrac{4}{5},\ \dfrac{2}{8},\ \dfrac{3}{10},\ \dfrac{7}{20}$
  \task $\dfrac{5}{2},\ \dfrac{2}{3},\ \dfrac{4}{7},\ \dfrac{5}{4}$
  \task $\dfrac{7}{12},\ \dfrac{2}{6},\ \dfrac{1}{4},\ \dfrac{3}{2}$
\end{tasks}

\ex{2} Marca en la recta les fraccions $\dfrac{1}{4},\ \dfrac{3}{4},\ \dfrac{5}{4},\ \dfrac{3}{2}$.\par\smallskip{\centering \begin{tikzpicture}[baseline=0]\draw[-latex,thick] (-0.3,0)--(9.50,0);\draw[muted] (1.125,0.1)--(1.125,-0.1);\draw[muted] (2.250,0.1)--(2.250,-0.1);\draw[muted] (3.375,0.1)--(3.375,-0.1);\draw[muted] (5.625,0.1)--(5.625,-0.1);\draw[muted] (6.750,0.1)--(6.750,-0.1);\draw[muted] (7.875,0.1)--(7.875,-0.1);\draw[thick] (0.00,0.18)--(0.00,-0.18) node[below=1pt,font=\small]{$0$};\draw[thick] (4.50,0.18)--(4.50,-0.18) node[below=1pt,font=\small]{$1$};\draw[thick] (9.00,0.18)--(9.00,-0.18) node[below=1pt,font=\small]{$2$};\end{tikzpicture}\par}

\ex{3} Escriu una fracció compresa entre $\dfrac{2}{5}$ i $\dfrac{3}{5}$, i una altra entre $\dfrac{1}{3}$ i $\dfrac{1}{2}$.

\ex{3} Tres amics llegeixen el mateix llibre: l'Anna n'ha llegit $\dfrac{3}{4}$, en Biel $\dfrac{5}{8}$ i la Clara $\dfrac{7}{10}$. Qui n'ha llegit més? Ordena'ls.


% =====================================================
\seccio{8}{Sumes i restes}

\ex{1} Calcula i simplifica el resultat:\begin{tasks}(3)
  \task $\dfrac{4}{7}+\dfrac{8}{7}$
  \task $\dfrac{2}{3}+\dfrac{1}{5}$
  \task $\dfrac{5}{2}-\dfrac{1}{8}$
  \task $\dfrac{3}{6}-\dfrac{4}{12}$
  \task $\dfrac{5}{8}+\dfrac{1}{5}$
  \task $\dfrac{14}{15}-\dfrac{2}{3}$
\end{tasks}\recorda{Si els denominadors són diferents, primer reduïm a \textbf{comú denominador}. \textbf{Mai} sumem els denominadors: $\dfrac{1}{2}+\dfrac{1}{3}=\dfrac{3}{6}+\dfrac{2}{6}=\dfrac{5}{6}$.}

\ex{1} Calcula i simplifica (els nombres naturals tenen denominador 1):\begin{tasks}(3)
  \task $\dfrac{3}{4}+\dfrac{6}{8}$
  \task $4+\dfrac{25}{9}$
  \task $\dfrac{25}{8}-2$
  \task $\dfrac{3}{5}+1$
  \task $5-\dfrac{3}{25}$
  \task $\dfrac{13}{6}+\dfrac{2}{7}$
\end{tasks}

\ex{2} Calcula i simplifica el resultat:\begin{tasks}(3)
  \task $\dfrac{12}{15}+\dfrac{8}{9}$
  \task $\dfrac{7}{12}-\dfrac{3}{10}$
  \task $\dfrac{7}{6}+\dfrac{8}{12}$
  \task $\dfrac{18}{5}-\dfrac{3}{10}$
  \task $\dfrac{5}{3}-\dfrac{9}{24}$
  \task $\dfrac{35}{36}+\dfrac{2}{9}$
\end{tasks}

\ex{2} Calcula i simplifica:\begin{tasks}(2)
  \task $\dfrac{4}{5}+\dfrac{2}{7}-\dfrac{1}{2}$
  \task $\dfrac{6}{8}+\dfrac{4}{3}-\dfrac{8}{12}$
  \task $\dfrac{3}{2}+\dfrac{6}{7}-\dfrac{2}{4}$
  \task $6-\dfrac{1}{2}-\dfrac{4}{12}$
  \task $\dfrac{2}{3}+\dfrac{1}{5}-\dfrac{1}{6}$
  \task $\dfrac{4}{3}+\dfrac{1}{6}-\dfrac{2}{7}$
\end{tasks}

\ex{3} Calcula i simplifica:\begin{tasks}(2)
  \task $\dfrac{2}{10}+\dfrac{5}{60}-\dfrac{1}{30}$
  \task $\dfrac{12}{5}+2-\dfrac{1}{25}$
  \task $\dfrac{6}{9}+\dfrac{1}{3}-\dfrac{2}{6}$
  \task $\dfrac{10}{3}-\dfrac{2}{6}-\dfrac{4}{5}+\dfrac{1}{2}$
  \task $4-\dfrac{1}{2}+\dfrac{1}{3}+2$
  \task $\dfrac{7}{9}+\dfrac{1}{3}-\dfrac{2}{5}$
\end{tasks}

\ex{3} Omplim una piscina en tres dies: el primer dia, $\dfrac{1}{4}$; el segon, $\dfrac{2}{5}$, i el tercer, $\dfrac{1}{6}$. Quina fracció de la piscina falta per omplir?


% =====================================================
\seccio{9}{Multiplicació i divisió}

\ex{1} Calcula i simplifica:\begin{tasks}(3)
  \task $\dfrac{6}{4}\cdot \dfrac{2}{5}$
  \task $\dfrac{3}{2}\cdot \dfrac{2}{3}$
  \task $2\cdot \dfrac{5}{6}$
  \task $\dfrac{3}{2}\cdot \dfrac{4}{7}\cdot 3$
  \task $\dfrac{3}{4}\cdot \dfrac{10}{2}\cdot \dfrac{1}{6}$
  \task $\dfrac{4}{3}\cdot \dfrac{2}{8}\cdot \dfrac{6}{2}$
\end{tasks}\recorda{$\dfrac{a}{b}\cdot\dfrac{c}{d}=\dfrac{a\cdot c}{b\cdot d}$. \qquad $\dfrac{a}{b}:\dfrac{c}{d}=\dfrac{a\cdot d}{b\cdot c}$ \ (productes creuats). \ La \textbf{inversa} de $\dfrac{a}{b}$ és $\dfrac{b}{a}$.}

\ex{1} Calcula (la fracció d'una fracció és un producte):\begin{tasks}(4)
  \task $\dfrac{2}{3}$ de $\dfrac{3}{5}$
  \task $\dfrac{4}{10}$ de $\dfrac{1}{2}$
  \task $\dfrac{4}{3}$ de $\dfrac{6}{2}$
  \task $\dfrac{3}{5}$ de $\dfrac{5}{7}$
\end{tasks}

\ex{1} Quines parelles són fraccions inverses? (El seu producte ha de ser 1.)\begin{tasks}(3)
  \task $\dfrac{2}{5},\ \dfrac{5}{3}$
  \task $\dfrac{4}{3},\ \dfrac{12}{16}$
  \task $\dfrac{7}{9},\ \dfrac{9}{7}$
  \task $\dfrac{2}{4},\ \dfrac{4}{2}$
  \task $\dfrac{5}{6},\ \dfrac{36}{30}$
  \task $\dfrac{12}{13},\ \dfrac{26}{12}$
\end{tasks}

\ex{1} Escriu la inversa de:\begin{tasks}(4)
  \task $\dfrac{4}{13}$
  \task $\dfrac{3}{7}$
  \task $\dfrac{15}{8}$
  \task $6$
\end{tasks}

\ex{2} Calcula i simplifica (amb divisions seguides, d'esquerra a dreta):\begin{tasks}(3)
  \task $\dfrac{4}{3}:\dfrac{5}{6}$
  \task $\dfrac{5}{4}:2$
  \task $\dfrac{4}{8}:\dfrac{6}{3}$
  \task $\dfrac{2}{6}:\dfrac{5}{4}:\dfrac{1}{3}$
  \task $\dfrac{3}{4}:\dfrac{12}{5}:\dfrac{2}{5}$
  \task $6:\dfrac{2}{3}:\dfrac{1}{5}$
\end{tasks}

\ex{2} Calcula i simplifica:\begin{tasks}(2)
  \task $\dfrac{3}{4}\cdot \dfrac{5}{8}:\dfrac{2}{5}$
  \task $\dfrac{5}{4}:\dfrac{1}{3}\cdot \dfrac{6}{5}$
  \task $\dfrac{5}{12}\cdot \dfrac{3}{2}:\dfrac{1}{6}$
  \task $\dfrac{6}{5}:\dfrac{7}{4}\cdot \dfrac{2}{10}$
\end{tasks}

\ex{3} Amb $12$ litres de suc, quantes ampolles de $\dfrac{3}{4}$ de litre podem omplir? I quina fracció cal multiplicar per $\dfrac{3}{4}$ per a obtindre $\dfrac{9}{10}$?


% =====================================================
\seccio{10}{Operacions combinades}

\ex{2} Calcula i simplifica:\begin{tasks}(2)
  \task $\dfrac{5}{2}+\dfrac{1}{3}\cdot \dfrac{4}{5}$
  \task $\dfrac{21}{4}-\dfrac{2}{5}:\dfrac{4}{3}$
  \task $\dfrac{4}{3}:\dfrac{2}{6}+\dfrac{1}{3}\cdot \dfrac{5}{4}$
  \task $\dfrac{4}{3}\cdot \left(\dfrac{2}{5}+\dfrac{1}{3}\right)$
  \task $\left(\dfrac{2}{5}+1\right):\dfrac{8}{3}$
  \task $\left(\dfrac{5}{4}+1\right)\cdot \dfrac{6}{7}$
\end{tasks}\recorda{Ordre: 1r \textbf{parèntesis}; 2n \textbf{multiplicacions i divisions} (d'esquerra a dreta); 3r \textbf{sumes i restes} (d'esquerra a dreta).}

\ex{2} Calcula i simplifica:\begin{tasks}(2)
  \task $\left(\dfrac{2}{5}+\dfrac{3}{7}\right)\cdot \dfrac{1}{2}-\dfrac{4}{10}$
  \task $\left(\dfrac{5}{4}+\dfrac{2}{3}+\dfrac{1}{12}\right)\cdot \dfrac{5}{4}-1$
  \task $\dfrac{5}{4}:\left(\dfrac{6}{4}-1\right)+\dfrac{2}{5}$
  \task $\dfrac{4}{3}-\left(\dfrac{1}{2}+\dfrac{2}{7}\right)+\dfrac{4}{5}\cdot \dfrac{2}{4}$
\end{tasks}

\ex{3} Calcula i simplifica:\begin{tasks}(2)
  \task $\left(\dfrac{1}{4}+\dfrac{2}{3}-\dfrac{1}{12}\right)\cdot 3+\left(\dfrac{3}{2}-\dfrac{5}{4}\right):\dfrac{1}{2}$
  \task $4\cdot \left(\dfrac{2}{3}+\dfrac{4}{6}-1\right)+\dfrac{2}{3}$
  \task $\left(\dfrac{3}{2}+\dfrac{1}{7}\right)\cdot \left(\dfrac{5}{3}-\dfrac{7}{15}\right)$
  \task $\dfrac{12}{10}:\dfrac{5}{4}+\dfrac{1}{2}\cdot \dfrac{13}{5}$
  \task $\left(\dfrac{4}{3}+\dfrac{2}{5}\cdot \dfrac{1}{3}\right):\dfrac{4}{12}$
  \task $\dfrac{3}{2}:\dfrac{1}{4}\cdot \dfrac{4}{20}+\dfrac{3}{2}\cdot \dfrac{7}{5}$
\end{tasks}

\ex{3} Calcula i simplifica:\begin{tasks}(1)
  \task $\left(\dfrac{5}{4}+\dfrac{2}{3}-1\right)\cdot \dfrac{3}{5}+\left(\dfrac{2}{5}+\dfrac{4}{3}\right)\cdot 2$
  \task $\dfrac{3}{8}+\dfrac{4}{5}\cdot \dfrac{1}{3}+\left(\dfrac{2}{6}-\dfrac{1}{4}\right)\cdot 3$
\end{tasks}


% =====================================================
\seccio{11}{Problemes amb fraccions}

\ex{1} D'un sopar, en Joan se n'ha menjat $\dfrac{1}{7}$, l'Elena $\dfrac{2}{5}$ i la Maite $\dfrac{1}{8}$.\begin{tasks}(1)
  \task Quant s'han menjat entre tots tres?
  \task Quant n'ha sobrat?
  \task Qui ha menjat més i qui menys?
\end{tasks}\recorda{Llig l'enunciat dues vegades, fes un \textbf{dibuix} (una barra que representa el total) i, al final, \textbf{comprova} que el resultat té sentit.}

\ex{1} D'un pastís de 24 porcions ens n'hem menjat $\dfrac{5}{8}$. Quantes porcions queden?

\ex{2} D'una pizza, en Carles se'n menja $\dfrac{3}{6}$, en Miquel $\dfrac{3}{21}$ i la Lorena $\dfrac{4}{14}$.\begin{tasks}(1)
  \task Quina fracció de la pizza ha sobrat?
  \task Arriba la Mireia amb una altra pizza igual. Com que els altres ja estan farts, ella es menja el que havia sobrat i $\dfrac{3}{5}$ de la seua pizza. Quanta pizza s'ha menjat ella?
\end{tasks}

\ex{2} En un poble de 5250 habitants, la meitat tenen més de 65 anys i $\dfrac{1}{5}$ en tenen menys de 18.\begin{tasks}(1)
  \task Quina fracció representen els que tenen entre 18 i 65 anys?
  \task Quants habitants hi ha en cada franja d'edat?
\end{tasks}

\ex{2} Un bidó de 60 litres està ple. Primer se'n gasten $\dfrac{2}{4}$ i després $\dfrac{1}{3}$ del que quedava. Finalment, el bidó bolca i es vessa la meitat del que hi ha. Quants litres queden al final?

\ex{3} Una família gasta $\dfrac{2}{6}$ del sou en la hipoteca, $\dfrac{1}{5}$ en menjar, $\dfrac{2}{7}$ en el cotxe i la resta en altres despeses. Quina fracció representen les altres despeses?

\ex{3} Una família ingressa $5400$\,€ al mes. Se'n gasta $\dfrac{2}{6}$ en la hipoteca, $\dfrac{1}{5}$ en menjar i $\dfrac{1}{9}$ en el cotxe.\begin{tasks}(1)
  \task Quants diners gasta en cada cosa i quants li'n queden per a altres despeses?
  \task Quina fracció del total representen els diners que sobren?
\end{tasks}

\ex{3} Un dipòsit està ple fins als $\dfrac{3}{4}$. Gastem $\dfrac{2}{3}$ del que hi ha i hi queden 50 litres. Quina capacitat té el dipòsit?

\vfill
\begin{center}\color{muted}\small\faCheckCircle[regular]\ \ Ànim! Si et costa algun exercici, repassa els apunts interactius de la unitat.\end{center}

\end{document}
